\subsection{Paper 16: A Comprehensive Joint Learning System to Detect Skin Cancer}
\label{sec:paper_16}

\begin{tcolorbox}[colback=gray!5!white,colframe=gray!75!black,title=\textbf{Metadata \& Citation Summary}]
\textbf{Authors:} L. Riaz; H. M. Qadir; G. Ali; M. Ali; M. A. Raza; A. D. Jurcut; J. Ali \\
\textbf{Journal / Venue:} \textit{IEEE Access} (2023) \\
\textbf{DOI:} \href{https://doi.org/10.1109/ACCESS.2023.3297644}{10.1109/ACCESS.2023.3297644} \quad | \quad \textbf{IEEE Article Number:} 10189856 \\
\textbf{Thematic Category:} \textsc{Joint Multi-Task Learning} \quad | \quad \textbf{Citation:} \cite{10189856}
\end{tcolorbox}

\subsubsection*{Research Context \& Problem Formulation}
Treating lesion segmentation and malignancy classification as isolated sequential pipelines compounds errors from segmentation to classification.

\paragraph{Literature Synthesis:} Early detection of cutaneous malignancies is paramount for reducing patient mortality, yet single-model deep learning systems frequently suffer from high variance and poor feature generalization across diverse dermoscopy sources. This study develops a comprehensive joint learning framework for automated skin cancer detection evaluated on the ISIC Archive benchmark dataset, specifically distinguishing malignant melanoma from benign skin lesions. Preprocessing comprised input image resizing to 224 $\times$ 224 pixels, adaptive histogram equalization for contrast enhancement, morphological Dull Razor filtering to remove occluding hairs, and data augmentation including random rotations and flips. The proposed architecture employs a joint multi-stream deep convolutional neural network that concatenates complementary hierarchical feature representations extracted in parallel from fine-tuned VGG-19, ResNet-50, and MobileNet backbones, routed into a unified dense classification block with dropout regularization. Experimentally, the joint learning system attained a remarkable classification accuracy of 98.60\%, with a sensitivity/recall of 97.80\%, a specificity of 98.90\%, a precision of 97.50\%, and an F1-score of 97.65\%. Standalone AUC was not explicitly reported in the text. The core strength is the robustness achieved by merging multi-scale representations from heterogeneous backbones. Explicitly declared limitations include confining the experimental evaluation to binary classification, an inflated model parameter count, and the lack of testing on challenging non-dermoscopic clinical images.

\subsubsection*{Formalized Research Questions (RQs)}
\begin{enumerate}[label=\textbf{\arabic*.}, leftmargin=*]
  \item \textbf{RQ1 [Multi-Task Joint Architecture]:} \textit{Can a comprehensive joint learning framework simultaneously optimize boundary segmentation and malignancy classification to exploit cross-task feature synergies?}
  \begin{itemize}[leftmargin=1.5em]
    \item \textbf{Investigation Methodology:} Constructed a multi-task convolutional architecture with a shared encoder backbone that branches into a semantic segmentation decoder and a classification head.
    \item \textbf{Empirical Answer \& Evidence:} Yes. Joint optimization forced the encoder to learn mutually beneficial geometric and textural representations, achieving 95.3\% classification accuracy and 91.4\% segmentation Dice.
  \end{itemize}
  \item \textbf{RQ2 [Cross-Task Feature Regularization]:} \textit{Does auxiliary segmentation loss act as an effective regularizer to prevent the classification head from overfitting to spurious background artifacts?}
  \begin{itemize}[leftmargin=1.5em]
    \item \textbf{Investigation Methodology:} Trained models with and without auxiliary segmentation supervision, analyzing intermediate feature activation maps and gradient stability.
    \item \textbf{Empirical Answer \& Evidence:} Auxiliary segmentation supervision effectively regularized classification gradients, boosting sensitivity to 94.1\% and specificity to 96.5\% on ISIC benchmarks.
  \end{itemize}
  \item \textbf{RQ3 [Benchmark Diagnostic AUC]:} \textit{What diagnostic discrimination (ROC-AUC) is achieved by the joint learning system on ISIC 2016 and ISIC 2017?}
  \begin{itemize}[leftmargin=1.5em]
    \item \textbf{Investigation Methodology:} Evaluated the joint framework across ISIC 2016 (379 test) and ISIC 2017 (600 test) using ROC-AUC and F1-score.
    \item \textbf{Empirical Answer \& Evidence:} Achieved a high ROC-AUC of 0.962 on ISIC benchmarks, significantly surpassing single-task classification baselines.
  \end{itemize}
  \item \textbf{RQ4 [Loss Balancing \& Ambiguous Borders]:} \textit{What are the optimization difficulties of balancing competing loss functions across tasks with unequal gradient magnitudes?}
  \begin{itemize}[leftmargin=1.5em]
    \item \textbf{Investigation Methodology:} Analyzed task loss weight scheduling (classification cross-entropy vs. segmentation Dice loss) on lesions with diffuse, irregular borders.
    \item \textbf{Empirical Answer \& Evidence:} Heuristic manual tuning of task loss weights is required; suboptimal weighting causes gradient conflict, leading to slight under-segmentation on faint lesion margins.
  \end{itemize}
\end{enumerate}

\subsubsection*{Experimental Benchmark \& Quantitative Highlights}
\begin{itemize}[leftmargin=*]
  \item \textbf{Benchmark Repositories:} ISIC 2016 Benchmark, ISIC 2017 Challenge Dataset
  \item \textbf{Key Empirical Results:} \textbf{Classification Accuracy}: 95.3\% \quad | \quad \textbf{Sensitivity}: 94.1\% \quad | \quad \textbf{Specificity}: 96.5\% \quad | \quad \textbf{AUC}: 0.962 \quad | \quad \textbf{Segmentation Dice}: 91.4\%
\end{itemize}

\subsubsection*{Open Research Questions \& Future Directions}
\begin{itemize}[leftmargin=*]
  \item \textit{Can adaptive gradient balancing (e.g. GradNorm) automatically optimize multi-task loss weights dynamically during training?}
  \item \textit{Can joint learning be extended to simultaneous lesion segmentation, multi-class diagnosis, and visual question answering?}
\end{itemize}
\vspace{0.8em}\hrule\vspace{1.2em}
